%% kapa: exam=2 round=22 year=2011 period=3 subject=compensation_law subject_ko=감정평가 및 보상법규 seq=1 number=1 points=50
%% source: Q-Net 감정평가사 기출문제 (cst00309) artlSeq=5040389 file=extracted/1008/5040389/2011년 제22회 감정평가사 2차시험 기출문제/3교시.hwp
%% conversion: kordoc 4.24.51 -> tools/gen2.py (auto); review=proofread-sonnet; flags=none
\begin{longquestion}{【문제 1】}
A군에 사는 甲은 국토의 계획 및 이용에 관한 법률에 따라 지정된 개발제한구역 내에 과수원을 경영하고 있다. 甲은 영농의 편의를 위해 동 과수원 토지 내에 작은 소로(小路)를 개설하고, 종종 이웃 주민의 통행에도 제공해 왔다. A군은 甲의 과수원 부지가 속한 일단의 토지에 폐기물처리장을 건설하고자 하는 乙을 폐기물관리법에 따라 폐기물처리장 건설사업자로 지정하면서 동 처리장건설사업실시계획을 승인하였다. 甲과 乙 간에 甲 토지에 대한 협의매수가 성립되지 않아 乙은 甲 토지에 대한 수용재결을 신청하고, 관할 지방토지수용위원회의 수용재결을 받았다. 동 수용재결에서는 “사실상의 사도(私道)의 부지는 인근 토지에 대한 평가액의 3분의 1 이내로 평가한다.”고 규정하고 있는 토지 등의 취득 및 보상에 관한 법률 시행규칙(이하 “토지보상법시행규칙”) 제26조 제1항 제2호의 규정에 따라, 甲의 토지를 인근 토지가에 비하여 3분의 1의 가격으로 평가하였다. 이 수용재결에 대하여 이의가 있는 甲은 적절한 권리구제 수단을 강구하고자 한다. 다음의 물음에 답하시오. (50점) \par
\medskip\noindent (1) 토지보상액에 대해 불복하고자 하는 甲의 행정쟁송상 권리구제수단을 설명하시오. (20점) \par
\medskip\noindent (2) 甲이 제기한 쟁송에서 피고 측은 甲의 토지에 대한 보상액이 낮게 평가된 것은 토지보상법시행규칙 제26조 제1항 제2호의 규정에 의한 것으로서 적법하다고 주장한다. 피고의 주장에 대해 법적으로 판단하시오. (15점) \par
\medskip\noindent (3) 甲은 토지보상법시행규칙 제26조 제1항 제2호의 규정은 헌법 제23조상의 재산권보장 및 정당보상원칙을 위배하여 위헌적인 것이라고 주장한다. 甲의 주장을 관철할 수 있는 법적 수단을 설명하시오. (15점) \par
\end{longquestion}
